\numoptsfalse
\yearpaper{2018}{Version A \textperiodcentered{} 60 Questions \textperiodcentered{} 60 Marks \textperiodcentered{} 80 Minutes}{}
\begin{multicols}{2}

\QQ{A particle shows distance-time curve as shown in the figure. The maximum instantaneous velocity of the particle is around the point
\fig{\gaxes{Time}{Distance}{\draw (0.2,0.2) .. controls (0.8,0.6) and (1.0,1.2) .. (1.3,1.45) .. controls (1.5,1.6) and (1.7,1.7) .. (1.9,1.72);
\fill (0.3,0.25) circle (1.4pt) node[right,font=\tiny]{P}; \fill (0.9,0.75) circle (1.4pt) node[right,font=\tiny]{Q}; \fill (1.35,1.48) circle (1.4pt) node[below right,font=\tiny]{R}; \fill (1.9,1.72) circle (1.4pt) node[below,font=\tiny]{S};}}\ntag{A}{14}{keph102}{2}{}{Class XI Part 1 \textperiodcentered{} Ch 2}\nsrc{None found in NEET/AIPMT/AIEEE/JEE papers, 2005--2026 (NCERT concept application)}}
{\opt{P}{S}{R}{Q}}

\QQ{Which of the following graphs correctly represents the variation of `$g$' on the Earth?\ntag{A}{133}{keph107}{7}{}{Class XI Part 1 \textperiodcentered{} Ch 7}\nsrc{None found in NEET/AIPMT/AIEEE/JEE papers, 2005--2026 (NCERT concept application)}}
{\optfig{\gaxes{$r$}{$g$}{\draw (0,0) -- (0.8,1.2) -- (1.7,0); \draw[densely dashed] (0.8,1.2) -- (0.8,0) node[below,font=\tiny]{R};}}
{\gaxes{$r$}{$g$}{\draw (0,0) -- (0.8,1.2) .. controls (1.0,0.6) and (1.3,0.45) .. (1.8,0.4); \draw[densely dashed] (0.8,1.2) -- (0.8,0) node[below,font=\tiny]{R};}}
{\gaxes{$r$}{$g$}{\draw (0,0) -- (0.8,1.2) -- (1.8,1.2); \draw[densely dashed] (0.8,1.2) -- (0.8,0) node[below,font=\tiny]{R};}}
{\gaxes{$r$}{$g$}{\draw (0,1.2) -- (0.8,1.2) -- (1.7,0); \draw[densely dashed] (0.8,1.2) -- (0.8,0) node[below,font=\tiny]{R};}}}

\QQ{A cup of tea cools from $65.5^\circ$C to $62.5^\circ$C in 1 minute in a room at $22.5^\circ$C. How long will it take to cool from $46.5^\circ$C to $40.5^\circ$C in the same room?\ntag{A}{220}{keph203}{19}{}{Class XI Part 2 \textperiodcentered{} Ch 10}\nsrc{None found in NEET/AIPMT/AIEEE/JEE papers, 2005--2026 (NCERT concept application)}}
{\opt{4 minutes}{2 minutes}{1 minute}{3 minutes}}

\QQ{The dimensions of the ratio of magnetic flux ($\phi$) and permeability ($\mu$) are\ntag{A}{7}{keph101}{7}{}{Class XI Part 1 \textperiodcentered{} Ch 1}\nsrc{None found in NEET/AIPMT/AIEEE/JEE papers, 2005--2026 (NCERT concept application)}}
{\opt{$[\mathrm{M^0L^1T^0A^1}]$}{$[\mathrm{M^0L^{-3}T^0A^1}]$}{$[\mathrm{M^0L^1T^1A^{-1}}]$}{$[\mathrm{M^0L^2T^0A^1}]$}}

\QQ{A mass `$m$' on the surface of the Earth is shifted to a target equal to the radius of the Earth. If `$R$' is the radius and `$M$' is the mass of the Earth, then work done in this process is\ntag{D}{134}{keph107}{8}{ \textperiodcentered{} Exemplar Q8.32}{Class XI Part 1 \textperiodcentered{} Ch 7}}
{\opt{$\dfrac{mgR}{2}$}{$mgR$}{$2mgR$}{$\dfrac{mgR}{4}$}}

\QQ{First overtone frequency of a closed pipe of length `$l_1$' is equal to the $2^{\text{nd}}$ harmonic frequency of an open pipe of length `$l_2$'. The ratio $\dfrac{l_1}{l_2} =$\ntag{A}{292}{keph207}{15}{}{Class XI Part 2 \textperiodcentered{} Ch 14}\nsrc{NEET 2016 (Phase II) -- overtone of an open pipe equal to an overtone of a closed pipe, ratio of lengths}}
{\opt{$\dfrac34$}{$\dfrac43$}{$\dfrac32$}{$\dfrac23$}}

\QQ{The resistance $R = \dfrac VI$ where $V = (100\pm5)\unit{V}$ and $I = (10\pm0.2)\unit{A}$. The percentage error in R is\ntag{A}{3}{keph101}{3}{}{Class XI Part 1 \textperiodcentered{} Ch 1}\nsrc{None found in NEET/AIPMT/AIEEE/JEE papers, 2005--2026. \textit{Note: error analysis is not in the 2026-27 NCERT textbook.}}}
{\opt{5.2\%}{4.8\%}{7\%}{3\%}}

\QQ{A block rests on a rough inclined plane making an angle of $30^\circ$ with the horizontal. The coefficient of static friction between the block and the plane is 0.8. If the frictional force on the block is $10\unit{N}$, the mass of the block is ($g = 10\unit{m\,s^{-2}}$)\ntag{A}{61}{keph104}{13}{}{Class XI Part 1 \textperiodcentered{} Ch 4}\nsrc{None found in NEET/AIPMT/AIEEE/JEE papers, 2005--2026 (NCERT concept application)}}
{\opt{$1\unit{kg}$}{$2\unit{kg}$}{$3\unit{kg}$}{$4\unit{kg}$}}

\QQ{Two particles of masses $m_1$ and $m_2$ have equal kinetic energies. The ratio of their momenta is\ntag{A}{74}{keph105}{4}{}{Class XI Part 1 \textperiodcentered{} Ch 5}\nsrc{None found in NEET/AIPMT/AIEEE/JEE papers, 2005--2026 (NCERT concept application)}}
{\opt{$m_1:m_2$}{$m_2:m_1$}{$\sqrt{m_1}:\sqrt{m_2}$}{$m_1^2:m_2^2$}}

\QQ{The pressure at the bottom of a liquid tank is \textbf{\textit{not}} proportional to the\ntag{D}{183}{keph202}{4}{}{Class XI Part 2 \textperiodcentered{} Ch 9}}
{\opt{Acceleration due to gravity}{Density of the liquid}{Height of the liquid}{Area of the liquid surface}}

\QQ{A Carnot engine takes 300 calories of heat from a source at $500\unit{K}$ and rejects 150 calories of heat to the sink. The temperature of the sink is\ntag{A}{238}{keph204}{13}{}{Class XI Part 2 \textperiodcentered{} Ch 11}\nsrc{None found in NEET/AIPMT/AIEEE/JEE papers, 2005--2026 (NCERT concept application)}}
{\opt{$125\unit{K}$}{$250\unit{K}$}{$750\unit{K}$}{$1000\unit{K}$}}

\QQ{Pressure of an ideal gas is increased by keeping temperature constant. The kinetic energy of molecules\ntag{A}{250}{keph205}{7}{}{Class XI Part 2 \textperiodcentered{} Ch 12}\nsrc{None found in NEET/AIPMT/AIEEE/JEE papers, 2005--2026 (NCERT concept application)}}
{\opt{Decreases}{Increases}{Remains same}{Increases or decreases depending on the nature of gas}}

\QQ{A man weighing $60\unit{kg}$ is in a lift moving down with an acceleration of $1.8\unit{m\,s^{-2}}$. The force exerted by the floor on him is\ntag{A}{53}{keph104}{5}{}{Class XI Part 1 \textperiodcentered{} Ch 4}\nsrc{None found in NEET/AIPMT/AIEEE/JEE papers, 2005--2026 (NCERT concept application)}}
{\opt{$588\unit{N}$}{$480\unit{N}$}{Zero}{$696\unit{N}$}}

\QQ{Moment of inertia of a body about two perpendicular axes X and Y in the plane of lamina are $20\unit{kg\,m^2}$ and $25\unit{kg\,m^2}$ respectively. Its moment of inertia about an axis perpendicular to the plane of the lamina and passing through the point of intersection of X and Y axes is\ntag{A}{115}{keph106}{24}{}{Class XI Part 1 \textperiodcentered{} Ch 6}\nsrc{None found in NEET/AIPMT/AIEEE/JEE papers, 2005--2026. \textit{Note: the perpendicular-axis theorem is not in the 2026-27 NCERT textbook.}}}
{\opt{$5\unit{kg\,m^2}$}{$45\unit{kg\,m^2}$}{$12.5\unit{kg\,m^2}$}{$500\unit{kg\,m^2}$}}

\QQ{Two wires A and B are stretched by the same load. If the area of cross-section of wire `A' is double that of `B', then the stress on `B' is\ntag{A}{169}{keph201}{3}{}{Class XI Part 2 \textperiodcentered{} Ch 8}\nsrc{None found in NEET/AIPMT/AIEEE/JEE papers, 2005--2026 (NCERT concept application)}}
{\opt{Equal to that on A}{Twice that on A}{Half that on A}{Four times that on A}}

\QQ{The magnitude of point charge due to which the electric field $30\unit{cm}$ away has the magnitude $2\unit{N\,C^{-1}}$ will be\ntag{A}{16}{leph101}{16}{}{Class XII Part 1 \textperiodcentered{} Ch 1}\nsrc{None found in NEET/AIPMT/AIEEE/JEE papers, 2005--2026 (NCERT concept application)}}
{\opt{$2\times10^{-11}\unit{C}$}{$3\times10^{-11}\unit{C}$}{$5\times10^{-11}\unit{C}$}{$9\times10^{-11}\unit{C}$}}

\QQ{A mass of $1\unit{kg}$ carrying a charge of $2\unit{C}$ is accelerated through a potential of $1\unit{V}$. The velocity acquired by it is\ntag{A}{47}{leph102}{3}{}{Class XII Part 1 \textperiodcentered{} Ch 2}\nsrc{None found in NEET/AIPMT/AIEEE/JEE papers, 2005--2026 (NCERT concept application)}}
{\opt{$\sqrt2\unit{m\,s^{-1}}$}{$2\unit{m\,s^{-1}}$}{$\dfrac1{\sqrt2}\unit{m\,s^{-1}}$}{$\dfrac12\unit{m\,s^{-1}}$}}

\QQ{The force of repulsion between two identical positive charges when kept with a separation `$r$' in air is `$F$'. Half the gap between the two charges is filled by a dielectric slab of dielectric constant $= 4$. Then the new force of repulsion between those two charges becomes\ntag{A}{8}{leph101}{8}{}{Class XII Part 1 \textperiodcentered{} Ch 1}\nsrc{None found in NEET/AIPMT/AIEEE/JEE papers, 2005--2026 (NCERT concept application)}}
{\opt{$\dfrac F3$}{$\dfrac F2$}{$\dfrac F4$}{$\dfrac{4F}9$}}

\QQ{For the arrangement of capacitors as shown in the circuit, the effective capacitance between the points A and B is (capacitance of each capacitor is $4\unit{\mu F}$)
\fig{\begin{circuitikz}[scale=0.6,transform shape,font=\Large]
\draw (0,1.5) node[circ]{} node[left]{A} -- (1,1.5);
\draw (1,3) to[C,l=$4\,\mu\mathrm F$] (1,1.5) to[C,l_=$4\,\mu\mathrm F$] (1,0);
\draw (3,3) to[C,l=$4\,\mu\mathrm F$] (3,1.5) to[C,l_=$4\,\mu\mathrm F$] (3,0);
\draw (1,3) -- (5,3) to[C,l=$4\,\mu\mathrm F$] (5,0) -- (1,0);
\draw (3,1.5) -- (4.2,1.5) node[circ]{} node[right]{B};
\end{circuitikz}}\ntag{A}{72}{leph102}{28}{}{Class XII Part 1 \textperiodcentered{} Ch 2}\nsrc{None found in NEET/AIPMT/AIEEE/JEE papers, 2005--2026 (NCERT concept application)}}
{\opt{$4\unit{\mu F}$}{$2\unit{\mu F}$}{$1\unit{\mu F}$}{$8\unit{\mu F}$}}

\QQ{The work done to move a charge on an equipotential surface is\ntag{D}{54}{leph102}{10}{}{Class XII Part 1 \textperiodcentered{} Ch 2}}
{\opt{Infinity}{Less than 1}{Greater than 1}{Zero}}

\QQ{Two capacitors of $3\unit{\mu F}$ and $6\unit{\mu F}$ are connected in series and a potential difference of $900\unit{V}$ is applied across the combination. They are then disconnected and reconnected in parallel. The potential difference across the combination is\ntag{A}{73}{leph102}{29}{}{Class XII Part 1 \textperiodcentered{} Ch 2}\nsrc{None found in NEET/AIPMT/AIEEE/JEE papers, 2005--2026 (NCERT concept application)}}
{\opt{Zero}{$100\unit{V}$}{$200\unit{V}$}{$400\unit{V}$}}

\QQ{Ohm's Law is applicable to\ntag{D}{89}{leph103}{9}{}{Class XII Part 1 \textperiodcentered{} Ch 3}}
{\opt{Diode}{Transistor}{Electrolyte}{Conductor}}

\QQ{If the last band on the carbon resistor is absent, then the tolerance is\ntag{A}{86}{leph103}{6}{}{Class XII Part 1 \textperiodcentered{} Ch 3}\nsrc{None found in NEET/AIPMT/AIEEE/JEE papers, 2005--2026. \textit{Note: the resistor colour code is not in the 2026-27 NCERT textbook.}}}
{\opt{5\%}{20\%}{10\%}{15\%}}

\QQ{The effective resistance between P and Q for the following network is
\fig{\begin{circuitikz}[scale=0.55,transform shape,font=\Large]
\draw (2.5,4.5) node[circ]{} node[above]{R} to[R,l_=$3\,\Omega$] (1.25,2.25) to[R,l_=$4\,\Omega$] (0,0) node[circ]{} node[below]{P};
\draw (2.5,4.5) to[R,l=$3\,\Omega$] (3.75,2.25) to[R,l=$5\,\Omega$] (5,0) node[circ]{} node[below]{Q};
\draw (1.25,2.25) to[R,l=$6\,\Omega$] (3.75,2.25);
\end{circuitikz}}\ntag{A}{97}{leph103}{17}{}{Class XII Part 1 \textperiodcentered{} Ch 3}\nsrc{None found in NEET/AIPMT/AIEEE/JEE papers, 2005--2026 (NCERT concept application)}}
{\opt{$\dfrac1{12}\unit{\Omega}$}{$21\unit{\Omega}$}{$12\unit{\Omega}$}{$\dfrac1{21}\unit{\Omega}$}}

\QQ{Five identical resistors each of resistance $R = 1500\unit{\Omega}$ are connected to a $300\unit{V}$ battery as shown in the circuit. The reading of the ideal ammeter A is
\fig{\begin{circuitikz}[scale=0.55,transform shape,font=\Large]
\draw (0,0) to[battery1,l=$300\,\mathrm V$,invert] (0,2.4) -- (7.6,2.4);
\foreach \x in {1.2,2.4,4.4,5.6,6.8} {\draw (\x,2.4) to[R,l=R] (\x,0); \fill (\x,2.4) circle (2.5pt);}
\draw (0,0) -- (3.0,0) to[rmeter,t=A] (3.8,0) -- (6.8,0);
\draw (1.2,0) -- (1.2,-0.5) -- (5.6,-0.5) -- (5.6,0);
\end{circuitikz}}\ntag{A}{97}{leph103}{17}{}{Class XII Part 1 \textperiodcentered{} Ch 3}\nsrc{None found in NEET/AIPMT/AIEEE/JEE papers, 2005--2026 (NCERT concept application)}}
{\opt{$\dfrac15\unit{A}$}{$\dfrac35\unit{A}$}{$\dfrac25\unit{A}$}{$\dfrac45\unit{A}$}}

\QQ{Two cells of internal resistances $r_1$ and $r_2$ and of same emf are connected in series, across a resistor of resistance R. If the terminal potential difference across the cells of internal resistance $r_1$ is zero, then the value of R is\ntag{A}{95}{leph103}{15}{}{Class XII Part 1 \textperiodcentered{} Ch 3}\nsrc{None found in NEET/AIPMT/AIEEE/JEE papers, 2005--2026 (NCERT concept application)}}
{\opt{$R = 2(r_1+r_2)$}{$R = r_2-r_1$}{$R = r_1-r_2$}{$R = 2(r_1-r_2)$}}

\QQ{The I-V graphs for two different electrical appliances P and Q are shown in the diagram. If $R_P$ and $R_Q$ be the resistances of the devices, then
\fig{\gaxes{$V$}{$I$}{\draw (0,0) -- (1.3,1.7) node[right,lbl]{Q}; \draw (0,0) -- (1.9,1.0) node[right,lbl]{P}; \node[lbl,below left] at (0,0) {O};}}\ntag{A}{84}{leph103}{4}{}{Class XII Part 1 \textperiodcentered{} Ch 3}\nsrc{None found in NEET/AIPMT/AIEEE/JEE papers, 2005--2026 (NCERT concept application)}}
{\opt{$R_P = R_Q$}{$R_P>R_Q$}{$R_P<R_Q$}{$R_P = \dfrac{R_Q}{2}$}}

\QQ{The correct Biot-Savart law in vector form is\ntag{D}{114}{leph104}{8}{}{Class XII Part 1 \textperiodcentered{} Ch 4}}
{\opt{$d\vec B = \dfrac{\mu_0}{4\pi}\dfrac{I(\vec{dl}\times\vec r)}{r^2}$}{$d\vec B = \dfrac{\mu_0}{4\pi}\dfrac{I(\vec{dl}\times\vec r)}{r^3}$}{$d\vec B = \dfrac{\mu_0}{4\pi}\dfrac{I\,\vec{dl}}{r^2}$}{$d\vec B = \dfrac{\mu_0}{4\pi}\cdot\dfrac{I\,\vec{dl}}{r^3}$}}

\QQ{An electron is moving in a circle of radius $r$ in a uniform magnetic field B. Suddenly the field is reduced to $\dfrac B2$. The radius of the circular path now becomes\ntag{A}{112}{leph104}{6}{}{Class XII Part 1 \textperiodcentered{} Ch 4}\nsrc{None found in NEET/AIPMT/AIEEE/JEE papers, 2005--2026 (NCERT concept application)}}
{\opt{$\dfrac r2$}{$2r$}{$\dfrac r4$}{$4r$}}

\QQ{A charge $q$ is accelerated through a potential difference V. It is then passed normally through a uniform magnetic field, where it moves in a circle of radius $r$. The potential difference required to move it in a circle of radius $2r$ is\ntag{A}{112}{leph104}{6}{}{Class XII Part 1 \textperiodcentered{} Ch 4}\nsrc{None found in NEET/AIPMT/AIEEE/JEE papers, 2005--2026 (NCERT concept application)}}
{\opt{2\,V}{4\,V}{1\,V}{3\,V}}

\QQ{A cyclotron's oscillator frequency is $10\unit{MHz}$ and the operating magnetic field is $0.66\unit{T}$. If the radius of its dees is $60\unit{cm}$, then the kinetic energy of the proton beam produced by the accelerator is\ntag{A}{112}{leph104}{6}{}{Class XII Part 1 \textperiodcentered{} Ch 4}\nsrc{None found in NEET/AIPMT/AIEEE/JEE papers, 2005--2026. \textit{Note: the cyclotron is not in the 2026-27 NCERT textbook.}}}
{\opt{$9\unit{MeV}$}{$10\unit{MeV}$}{$7\unit{MeV}$}{$11\unit{MeV}$}}

\QQ{Needles $N_1$, $N_2$ and $N_3$ are made of a ferromagnetic, a paramagnetic and a diamagnetic substance respectively. A magnet when brought close to them will\ntag{D}{147}{leph105}{12}{}{Class XII Part 1 \textperiodcentered{} Ch 5}}
{\opt{Attract all three of them}{Attract $N_1$ strongly, $N_2$ weakly and repel $N_3$ weakly}{Attract $N_1$ strongly but repel $N_2$ and $N_3$ weakly}{Attract $N_1$ and $N_2$ strongly but repel $N_3$}}

\QQ{The strength of the Earth's magnetic field is\ntag{A}{136}{leph105}{1}{}{Class XII Part 1 \textperiodcentered{} Ch 5}\nsrc{None found in NEET/AIPMT/AIEEE/JEE papers, 2005--2026. \textit{Note: Earth's magnetism is not in the 2026-27 NCERT textbook.}}}
{\opt{Constant everywhere}{Zero everywhere}{Having very high value}{Varying from place to place on the Earth's surface}}

\QQ{A jet plane having a wing-span of $25\unit{m}$ is travelling horizontally towards east with a speed of $3600\unit{km/hour}$. If the Earth's magnetic field at the location is $4\times10^{-4}\unit{T}$ and the angle of dip is $30^\circ$, then, the potential difference between the ends of the wing is\ntag{A}{163}{leph106}{10}{}{Class XII Part 1 \textperiodcentered{} Ch 6}\nsrc{None found in NEET/AIPMT/AIEEE/JEE papers, 2005--2026. \textit{Note: Earth's magnetism (angle of dip) is not in the 2026-27 NCERT textbook.}}}
{\opt{$4\unit{V}$}{$5\unit{V}$}{$2\unit{V}$}{$2.5\unit{V}$}}

\QQ{Which of the following, represents the variation of inductive reactance ($X_L$) with the frequency of voltage source ($\nu$)?\ntag{A}{182}{leph107}{6}{}{Class XII Part 1 \textperiodcentered{} Ch 7}\nsrc{None found in NEET/AIPMT/AIEEE/JEE papers, 2005--2026 (NCERT concept application)}}
{\optfig{\gaxes{$\nu$}{$X_L$}{\draw (0,0) -- (1.5,1.5);}}{\gaxes{$\nu$}{$X_L$}{\draw (0.1,0.1) .. controls (0.9,0.15) and (1.4,0.6) .. (1.6,1.5);}}{\gaxes{$\nu$}{$X_L$}{\draw (0.1,0.1) .. controls (0.2,1.0) and (0.7,1.4) .. (1.6,1.5);}}{\gaxes{$\nu$}{$X_L$}{\draw (0.15,1.6) .. controls (0.3,0.4) and (0.8,0.2) .. (1.7,0.15);}}}

\QQ{The magnetic flux linked with a coil varies as $\phi = 3t^2+4t+9$. The magnitude of the emf induced at $t = 2$ seconds is\ntag{A}{157}{leph106}{4}{}{Class XII Part 1 \textperiodcentered{} Ch 6}\nsrc{None found in NEET/AIPMT/AIEEE/JEE papers, 2005--2026 (NCERT concept application)}}
{\opt{$8\unit{V}$}{$16\unit{V}$}{$32\unit{V}$}{$64\unit{V}$}}

\QQ{A $100\unit{W}$ bulb is connected to an AC source of $220\unit{V}$, $50\unit{Hz}$. Then the current flowing through the bulb is\ntag{A}{180}{leph107}{4}{}{Class XII Part 1 \textperiodcentered{} Ch 7}\nsrc{None found in NEET/AIPMT/AIEEE/JEE papers, 2005--2026 (NCERT concept application)}}
{\opt{$\dfrac5{11}\unit{A}$}{$\dfrac12\unit{A}$}{$2\unit{A}$}{$\dfrac34\unit{A}$}}

\QQ{In the series LCR circuit, the power dissipation is through\ntag{D}{191}{leph107}{15}{}{Class XII Part 1 \textperiodcentered{} Ch 7}}
{\opt{R}{L}{C}{Both L and C}}

\QQ{In Karnataka, the normal domestic power supply AC is $220\unit{V}$, $50\unit{Hz}$. Here $220\unit{V}$ and $50\unit{Hz}$ refer to\ntag{D}{180}{leph107}{4}{}{Class XII Part 1 \textperiodcentered{} Ch 7}}
{\opt{Peak value of voltage and frequency}{rms value of voltage and frequency}{Mean value of voltage and frequency}{Peak value of voltage and angular frequency}}

\QQ{A step-up transformer operates on a $230\unit{V}$ line and a load current of $2\unit{A}$. The ratio of primary and secondary windings is 1:25. Then the current in the primary is\ntag{A}{195}{leph107}{19}{}{Class XII Part 1 \textperiodcentered{} Ch 7}\nsrc{None found in NEET/AIPMT/AIEEE/JEE papers, 2005--2026 (NCERT concept application)}}
{\opt{$25\unit{A}$}{$50\unit{A}$}{$15\unit{A}$}{$12.5\unit{A}$}}

\QQ{The number of photons falling per second on a completely darkened plate to produce a force of $6.62\times10^{-5}\unit{N}$ is `$n$'. If the wavelength of the light falling is $5\times10^{-7}\unit{m}$, then $n =$ \rule{1cm}{0.4pt} $\times10^{22}$. ($h = 6.62\times10^{-34}\unit{J\,s}$)\ntag{A}{282}{leph203}{9}{}{Class XII Part 2 \textperiodcentered{} Ch 11}\nsrc{None found in NEET/AIPMT/AIEEE/JEE papers, 2005--2026 (NCERT concept application)}}
{\opt{1}{5}{0.2}{3.3}}

\QQ{An object is placed at the principal focus of a convex mirror. The image will be at\ntag{A}{224}{leph201}{4}{}{Class XII Part 2 \textperiodcentered{} Ch 9}\nsrc{None found in NEET/AIPMT/AIEEE/JEE papers, 2005--2026 (NCERT concept application)}}
{\opt{Centre of curvature}{Principal focus}{Infinity}{No image will be formed}}

\QQ{An object is placed at a distance of $20\unit{cm}$ from the pole of a concave mirror of focal length $10\unit{cm}$. The distance of the image formed is\ntag{A}{225}{leph201}{5}{}{Class XII Part 2 \textperiodcentered{} Ch 9}\nsrc{None found in NEET/AIPMT/AIEEE/JEE papers, 2005--2026 (NCERT concept application)}}
{\opt{$+20\unit{cm}$}{$+10\unit{cm}$}{$-20\unit{cm}$}{$-10\unit{cm}$}}

\QQ{A candle placed $25\unit{cm}$ from a lens forms an image on a screen placed $75\unit{cm}$ on the other side of the lens. The focal length and type of the lens should be\ntag{A}{235}{leph201}{15}{}{Class XII Part 2 \textperiodcentered{} Ch 9}\nsrc{None found in NEET/AIPMT/AIEEE/JEE papers, 2005--2026 (NCERT concept application)}}
{\opt{$+18.75\unit{cm}$ and convex lens}{$-18.75\unit{cm}$ and concave lens}{$+20.25\unit{cm}$ and convex lens}{$-20.25\unit{cm}$ and concave lens}}

\QQ{A plane wavefront of wavelength $\lambda$ is incident on a single slit of width $a$. The angular width of principal maximum is\ntag{A}{267}{leph202}{13}{}{Class XII Part 2 \textperiodcentered{} Ch 10}\nsrc{None found in NEET/AIPMT/AIEEE/JEE papers, 2005--2026 (NCERT concept application)}}
{\opt{$\dfrac\lambda a$}{$\dfrac{2\lambda}a$}{$\dfrac a\lambda$}{$\dfrac a{2\lambda}$}}

\QQ{In a Fraunhofer diffraction at a single slit, if yellow light illuminating the slit is replaced by blue light, then diffraction bands\ntag{A}{267}{leph202}{13}{}{Class XII Part 2 \textperiodcentered{} Ch 10}\nsrc{None found in NEET/AIPMT/AIEEE/JEE papers, 2005--2026 (NCERT concept application)}}
{\opt{Remain unchanged}{Become wider}{Disappear}{Become narrower}}

\QQ{In Young's double slit experiment, two wavelengths $\lambda_1 = 780\unit{nm}$ and $\lambda_2 = 520\unit{nm}$ are used to obtain interference fringes. If the $n^{\text{th}}$ bright band due to $\lambda_1$ coincides with $(n+1)^{\text{th}}$ bright band due to $\lambda_2$, then the value of $n$ is\ntag{A}{273}{leph202}{19}{}{Class XII Part 2 \textperiodcentered{} Ch 10}\nsrc{None found in NEET/AIPMT/AIEEE/JEE papers, 2005--2026 (NCERT concept application)}}
{\opt{4}{3}{2}{6}}

\QQ{In Young's double slit experiment, slits are separated by $2\unit{mm}$ and the screen is placed at a distance of $1.2\unit{m}$ from the slits. Light consisting of two wavelengths $6500\unit{\text{\AA}}$ and $5200\unit{\text{\AA}}$ are used to obtain interference fringes. Then the separation between the fourth bright fringes of two different patterns produced by the two wavelengths is\ntag{D}{273}{leph202}{19}{}{Class XII Part 2 \textperiodcentered{} Ch 10}}
{\opt{$0.312\unit{mm}$}{$0.123\unit{mm}$}{$0.213\unit{mm}$}{$0.412\unit{mm}$}}

\QQ{The maximum kinetic energy of emitted photoelectrons depends on\ntag{D}{280}{leph203}{7}{}{Class XII Part 2 \textperiodcentered{} Ch 11}}
{\opt{Intensity of incident radiation}{Frequency of incident radiation}{Speed of incident radiation}{Number of photons in the incident radiation}}

\QQ{A proton and an $\alpha$ particle are accelerated through the same potential difference V. The ratio of their de Broglie wavelengths is\ntag{A}{285}{leph203}{12}{}{Class XII Part 2 \textperiodcentered{} Ch 11}\nsrc{None found in NEET/AIPMT/AIEEE/JEE papers, 2005--2026 (NCERT concept application)}}
{\opt{$\sqrt2$}{$2\sqrt2$}{$\sqrt3$}{$2\sqrt3$}}

\QQ{The total energy of an electron revolving in the second orbit of hydrogen atom is\ntag{D}{300}{leph204}{11}{}{Class XII Part 2 \textperiodcentered{} Ch 12}}
{\opt{$-13.6\unit{eV}$}{$-1.51\unit{eV}$}{$-3.4\unit{eV}$}{Zero}}

\QQ{The period of revolution of an electron in the ground state of hydrogen atom is T. The period of revolution of the electron in the first excited state is\ntag{A}{299}{leph204}{10}{}{Class XII Part 2 \textperiodcentered{} Ch 12}\nsrc{None found in NEET/AIPMT/AIEEE/JEE papers, 2005--2026 (NCERT concept application)}}
{\opt{2T}{4T}{6T}{8T}}

\QQ{The energy equivalent to a substance of mass $1\unit{g}$ is\ntag{A}{310}{leph205}{5}{}{Class XII Part 2 \textperiodcentered{} Ch 13}\nsrc{None found in NEET/AIPMT/AIEEE/JEE papers, 2005--2026 (NCERT concept application)}}
{\opt{$18\times10^{13}\unit{J}$}{$9\times10^{13}\unit{J}$}{$18\times10^6\unit{J}$}{$9\times10^6\unit{J}$}}

\QQ{The half-life of tritium is 12.5 years. What mass of tritium of initial mass $64\unit{mg}$ will remain undecayed after 50 years?\ntag{A}{314}{leph205}{9}{}{Class XII Part 2 \textperiodcentered{} Ch 13}\nsrc{None found in NEET/AIPMT/AIEEE/JEE papers, 2005--2026. \textit{Note: the radioactive decay law (half-life) is not in the 2026-27 NCERT textbook.}}}
{\opt{$32\unit{mg}$}{$8\unit{mg}$}{$16\unit{mg}$}{$4\unit{mg}$}}

\QQ{In a CE amplifier, the input ac signal to be amplified is applied across\ntag{A}{323}{leph206}{1}{}{Class XII Part 2 \textperiodcentered{} Ch 14}\nsrc{None found in NEET/AIPMT/AIEEE/JEE papers, 2005--2026. \textit{Note: the transistor amplifier is not in the 2026-27 NCERT textbook.}}}
{\opt{Forward biased emitter-base junction}{Reverse biased collector-base junction}{Reverse biased emitter-base junction}{Forward biased collector-base junction}}

\QQ{If $A = 1$ and $B = 0$, then in terms of Boolean algebra, $A+\bar B =$\ntag{A}{323}{leph206}{1}{}{Class XII Part 2 \textperiodcentered{} Ch 14}\nsrc{None found in NEET/AIPMT/AIEEE/JEE papers, 2005--2026. \textit{Note: Boolean algebra (logic gates) is not in the 2026-27 NCERT textbook.}}}
{\opt{B}{$\bar B$}{A}{$\bar A$}}

\QQ{The density of an electron-hole pair in a pure germanium is $3\times10^{16}\unit{m^{-3}}$ at room temperature. On doping with aluminium, the hole density increases to $4.5\times10^{22}\unit{m^{-3}}$. Now the electron density (in $\mathrm{m^{-3}}$) in doped germanium will be\ntag{A}{332}{leph206}{10}{}{Class XII Part 2 \textperiodcentered{} Ch 14}\nsrc{AIPMT 2011 (Mains) -- pure semiconductor doped with a trivalent impurity, new electron density from $n_e n_h = n_i^2$}}
{\opt{$1\times10^{10}$}{$2\times10^{10}$}{$0.5\times10^{10}$}{$4\times10^{10}$}}

\QQ{The dc common emitter current gain of a n-p-n transistor is 50. The potential difference applied across the collector and emitter of a transistor used in CE configuration is, $V_{CE} = 2\unit{V}$. If the collector resistance, $R_C = 4\unit{k\Omega}$, the base current ($I_B$) and the collector current ($I_C$) are\ntag{A}{323}{leph206}{1}{}{Class XII Part 2 \textperiodcentered{} Ch 14}\nsrc{None found in NEET/AIPMT/AIEEE/JEE papers, 2005--2026. \textit{Note: the transistor is not in the 2026-27 NCERT textbook.}}}
{\opt{$I_B = 10\unit{\mu A}$, $I_C = 0.5\unit{mA}$}{$I_B = 0.5\unit{\mu A}$, $I_C = 10\unit{mA}$}{$I_B = 5\unit{\mu A}$, $I_C = 1\unit{mA}$}{$I_B = 1\unit{\mu A}$, $I_C = 0.5\unit{mA}$}}

\QQ{The radius of the Earth is $6400\unit{km}$. If the height of an antenna is $500\unit{m}$, then its range is\ntag{A}{209}{leph108}{9}{}{Class XII Part 1 \textperiodcentered{} Ch 8}\nsrc{AIEEE 2012 -- same question: antenna of height 500 m, $R=6.4\times10^6$ m, maximum range. \textit{Note: antenna range (communication systems) is not in the 2026-27 NCERT textbook.}}}
{\opt{$800\unit{km}$}{$100\unit{km}$}{$80\unit{km}$}{$10\unit{km}$}}

\QQ{A space station is at a height equal to the radius of the Earth. If `$v_E$' is the escape velocity on the surface of the Earth, the same on the space station is \rule{1cm}{0.4pt} times $v_E$.\ntag{A}{135}{keph107}{9}{}{Class XI Part 1 \textperiodcentered{} Ch 7}\nsrc{None found in NEET/AIPMT/AIEEE/JEE papers, 2005--2026 (NCERT concept application)}}
{\opt{$\dfrac12$}{$\dfrac14$}{$\dfrac1{\sqrt2}$}{$\dfrac1{\sqrt3}$}}

%%BEGIN-ANALYSIS
\end{multicols}
\clearpage
\begin{center}{\pgfornament[width=5cm,color=gilt]{87}}\par
{\fontsize{20}{24}\selectfont\color{navy}\scshape KCET 2018 \textperiodcentered{} Paper Analysis}\par
{\color{gilt}\pgfornament[width=5cm,color=gilt,symmetry=h]{87}}\end{center}
\vspace{4pt}
\begin{multicols}{2}
\textbf{\color{crimson}1. Direct vs Applied}\par\smallskip
\begin{tabular}{@{}lr@{}}\toprule
Direct from NCERT textbook (2026-27) & 10\\
Direct from NCERT Exemplar & 1\\
Applied (concept in NCERT) & 49\\\midrule
\textbf{Total} & \textbf{60}\\\bottomrule\end{tabular}\par\smallskip
{\footnotesize Direct: Q 10, 20, 22, 28, 32, 38, 39, 48, 49, 51.\\ Exemplar: Q 5.}\par\medskip
\textbf{\color{crimson}2. Where the applied questions resemble earlier exams}\par\smallskip
Of the 49 applied questions, \textbf{3} closely resemble a question from an earlier national exam:
NEET/AIPMT/AIIMS 2, JEE Main/Advanced 0, AIEEE 1.
The other 46 (Q 1, 2, 3, 4, 7, 8, 9, 11, 12, 13, 14, 15, 16, 17, 18, 19, 21, 23, 24, 25, 26, 27, 29, 30, 31, 33, 34, 35, 36, 37, 40, 41, 42, 43, 44, 45, 46, 47, 50, 52, 53, 54, 55, 56, 58, 60) had no close match in the NEET, AIPMT, AIIMS, AIEEE and JEE papers of 2005--2026 that were searched; they are straightforward applications of NCERT concepts. \textbf{11} questions (Q 7, 14, 23, 31, 33, 34, 54, 55, 56, 58, 59) test topics no longer in the 2026-27 NCERT textbook.\par
{\footnotesize Strongest repeats: Q59 (AIEEE 2012, identical), Q57 (AIPMT 2011), Q6 (NEET 2016).}\par\medskip
\textbf{\color{crimson}3. Difficulty level}\par\smallskip
\begin{tabular}{@{}lr@{}}\toprule
Easy & 43\\ Moderate & 17\\ Hard & 0\\\bottomrule\end{tabular}\par\smallskip
{\footnotesize Overall: \textbf{easy}; mostly one-step formula and statement questions, with no question needing more than two steps. Class XI: 16 questions, Class XII: 44. (Difficulty is the editor's assessment.)}\par\medskip
\columnbreak
\textbf{\color{crimson}4. Chapter-wise distribution}\par\smallskip
{\footnotesize\begin{tabular}{@{}p{3.6cm}c@{\,}r@{\,}r@{\,}r@{}}\toprule
Chapter & Cl. & Qs & Dir & App\\\midrule
Units \& Measurement & XI & 2 & 0 & 2\\
Motion in a Straight Line & XI & 1 & 0 & 1\\
Laws of Motion & XI & 2 & 0 & 2\\
Work, Energy \& Power & XI & 1 & 0 & 1\\
Rotational Motion & XI & 1 & 0 & 1\\
Gravitation & XI & 3 & 1 & 2\\
Mech.\ Prop.\ of Solids & XI & 1 & 0 & 1\\
Mech.\ Prop.\ of Fluids & XI & 1 & 1 & 0\\
Thermal Prop.\ of Matter & XI & 1 & 0 & 1\\
Thermodynamics & XI & 1 & 0 & 1\\
Kinetic Theory & XI & 1 & 0 & 1\\
Waves & XI & 1 & 0 & 1\\
Electric Charges \& Fields & XII & 2 & 0 & 2\\
Electrostatic Potential \& Capacitance & XII & 4 & 1 & 3\\
Current Electricity & XII & 6 & 1 & 5\\
Moving Charges \& Magnetism & XII & 4 & 1 & 3\\
Magnetism \& Matter & XII & 2 & 1 & 1\\
Electromagnetic Induction & XII & 2 & 0 & 2\\
Alternating Current & XII & 5 & 2 & 3\\
Electromagnetic Waves & XII & 1 & 0 & 1\\
Ray Optics & XII & 3 & 0 & 3\\
Wave Optics & XII & 4 & 1 & 3\\
Dual Nature & XII & 3 & 1 & 2\\
Atoms & XII & 2 & 1 & 1\\
Nuclei & XII & 2 & 0 & 2\\
Semiconductor Electronics & XII & 4 & 0 & 4\\
\bottomrule\end{tabular}}\par\medskip
\textbf{\color{crimson}5. Teaching takeaways}\par\smallskip
{\footnotesize\begin{itemize}[leftmargin=1em,nosep]
\item NCERT exercises and statements asked directly: Exercise 10.6 on two wavelengths (Q48, p.\,273), Ohm's-law limitations (p.\,89), power only in $R$ and rms 220 V (p.\,191, p.\,180), $E_2=-3.4$ eV (p.\,300); Q5 is Exemplar XI 8.32.
\item Eleven questions are on topics removed from the rationalised NCERT (errors, perpendicular-axis theorem, colour code, cyclotron, Earth's magnetism, half-life, transistor, Boolean algebra, antenna range).
\item Q59 repeats AIEEE 2012 word for word; old AIEEE and AIPMT papers are useful revision.
\item Q7 (errors) and Q23 (colour code) are lab-manual questions.
\end{itemize}}
\end{multicols}
%%END-ANALYSIS

